\section{Sharp stop models and route distinctions}
\label{sec:tpc67-countermodels}

The determinant identity separates exact rank information from the
strength of a particular quantitative certificate.  The following
models show that these stops must not be conflated.

\subsection{Exponential determinant loss with a uniform spectral gap}

\begin{example}[The path Gram]
\label{ex:tpc67-path-determinant-decay}
Let \(P_r\) be the path on \(r\) vertices, let \(A(P_r)\) be its adjacency
matrix, and fix
\[
 0<a<\frac12.
\]
Define the unit-diagonal matrix
\begin{equation}
 K_r:=I_r+aA(P_r).
 \label{eq:tpc67-path-Gram}
\end{equation}
Its eigenvalues are
\begin{equation}
 1+2a\cos\frac{k\pi}{r+1},
 \qquad k=1,\ldots,r.
 \label{eq:tpc67-path-spectrum}
\end{equation}
Consequently
\begin{equation}
 \lambda_{\min}(K_r)
 =
 1-2a\cos\frac{\pi}{r+1}
 >
 1-2a>0
 \label{eq:tpc67-path-uniform-gap}
\end{equation}
uniformly in \(r\).  In particular, \(K_r\) is a Gram matrix; one may
take \(B_r=K_r^{1/2}\).

Put \(D_r:=\det K_r\), with \(D_0:=1\).  Expansion along an endpoint of
the path gives
\begin{equation}
 D_0=D_1=1,
 \qquad
 D_r=D_{r-1}-a^2D_{r-2}
 \quad(r\ge2).
 \label{eq:tpc67-path-determinant-recurrence}
\end{equation}
If
\begin{equation}
 q_\pm:=\frac{1\pm\sqrt{1-4a^2}}2,
 \label{eq:tpc67-path-characteristic-roots}
\end{equation}
then
\begin{equation}
 D_r
 =
 \frac{q_+^{\,r+1}-q_-^{\,r+1}}{q_+-q_-}
 \asymp_a q_+^r,
 \qquad
 0<q_-<q_+<1.
 \label{eq:tpc67-path-determinant-decay}
\end{equation}
Thus the exact determinant mass of the square synthesis
\(B_r=K_r^{1/2}\) decays exponentially, while its smallest Gram
eigenvalue stays uniformly positive.
\end{example}

\begin{proof}
The path-adjacency spectrum is
\[
 \left\{2\cos\frac{k\pi}{r+1}:1\le k\le r\right\},
\]
which proves
\eqref{eq:tpc67-path-spectrum}--\eqref{eq:tpc67-path-uniform-gap}.
The tridiagonal determinant expansion gives
\eqref{eq:tpc67-path-determinant-recurrence}.  Solving its characteristic
equation \(q^2-q+a^2=0\) gives the displayed formula.  Since
\(0<q_-/q_+<1\), the factor
\[
 \frac{q_+\bigl(1-(q_-/q_+)^{r+1}\bigr)}{q_+-q_-}
\]
is bounded above and below by positive constants depending only on
\(a\), which proves the final comparison.
\end{proof}

\begin{remark}[A determinant-route stop is not a spectral stop]
\label{rem:tpc67-determinant-not-spectral-stop}
For the path family,
\[
 c_r\det K_r\asymp_a q_+^r
\]
is exponentially small, although
\(\lambda_{\min}(K_r)>1-2a\).  Therefore failure to obtain a
dimension-uniform lower bound through determinant mass does not prove
that the actual lower frame is small.  It only stops the determinant
conversion as a sufficiently strong route; direct spectral-radius,
Schur, forest, or other Gram estimates may still succeed.
\end{remark}

\subsection{Hall stops and activity-certificate stops}

\begin{example}[Support Hall deficiency is a genuine kernel]
\label{ex:tpc67-Hall-deficiency-kernel}
Let
\[
 B=\begin{pmatrix}1&1\end{pmatrix}.
\]
Both unit columns have the same unique output neighbor.  For the set of
both columns, the support neighborhood has size one, so Hall deficiency
is one.  Moreover,
\[
 K=B^*B=
 \begin{pmatrix}1&1\\1&1\end{pmatrix},
 \qquad
 \det K=0,
 \qquad
 \lambda_{\min}(K)=0.
\]
Here the support obstruction forces an actual kernel; this is a genuine
spectral stop.
\end{example}

\begin{example}[Positive cycle activity is only certificate failure]
\label{ex:tpc67-activity-certificate-failure}
For \(|\omega|=1\), let
\begin{equation}
 B_\omega
 :=
 \frac1{\sqrt2}
 \begin{pmatrix}
  1&1\\
  1&\omega
 \end{pmatrix}.
 \label{eq:tpc67-two-matching-cycle}
\end{equation}
The support is the complete bipartite graph \(K_{2,2}\).  It satisfies
all Hall inequalities, has two perfect matchings, and its square support
contains one alternating four-cycle.  Thus a unique-matching, forced-
matching certificate is unavailable.  Relative to either perfect
matching, the nonidentity cycle term has magnitude one, so
\(\mathcal A_\phi=1\) and the strict activity certificate
\(\mathcal A_\phi<1\) is also unavailable.

Nevertheless,
\begin{equation}
 \det B_\omega=\frac{\omega-1}{2}.
 \label{eq:tpc67-two-matching-determinant}
\end{equation}
For \(\omega=1\), matching terms cancel and the Gram has a kernel.  For
\(\omega=-1\),
\[
 B_{-1}^*B_{-1}=I_2,
 \qquad
 |\det B_{-1}|^2=1,
 \qquad
 \lambda_{\min}(B_{-1}^*B_{-1})=1.
\]
Hence activity at least one does not imply a kernel or even a small
gap.  It says only that support alone no longer certifies the literal
determinant without controlling the alternating-cycle holonomy.
\end{example}

\begin{remark}[Four distinct stop statements]
\label{rem:tpc67-stop-distinctions}
The examples and the exact determinant identity give four logically
different conclusions.
\begin{enumerate}[label=\textup{(\roman*)},leftmargin=3em]
 \item A support Hall deficiency forces
       \(\mathfrak D_X=\beta_X=0\).  This is a support-level kernel stop.
 \item Hall may hold while literal matching terms cancel, as for
       \(B_1\).  This is an actual coefficient kernel, not a support
       deficiency.
 \item Activity \(\mathcal A_\phi\ge1\), or failure of a unique-matching
       certificate, does not imply
       \(\mathfrak D_X=0\).  The matrix \(B_{-1}\) has no such certificate
       but has determinant mass and spectral gap equal to one.
 \item Even an exactly known positive determinant mass may be
       exponentially small while the true spectral gap is uniform, as
       in \cref{ex:tpc67-path-determinant-decay}.  This stops only the
       determinant-to-gap route at the desired scale.
\end{enumerate}
In particular, an empty certified-minor family gives a zero
\emph{certified lower sum}; it does not say that the omitted minors, the
full determinant mass, or the lower frame vanish.
\end{remark}

\begin{remark}[Claim level]
\label{rem:tpc67-countermodel-claim-level}
These are L0 finite-dimensional countermodels.  They separate route
diagnostics before any L1 arithmetic specialization.  None asserts that
the literal fixed-\(h_0\) collision core realizes the path family, the
two-matching phase, or a Hall defect.  Establishing any such property
for the physical coefficients is a separate arithmetic task.
\end{remark}
